\documentclass{article}
\usepackage{amsmath}
\usepackage{amsthm}
\usepackage[hidelinks]{hyperref}
\usepackage{amssymb}
\usepackage[
    left=2.15cm,
    right=2.15cm,
    top=1.90cm,
    bottom=1.90cm
]{geometry}


\title{Mathematical Tools for Economists}
\author{Leroy Zahno}
\date{Autumn Semester 2026}

\newtheorem{theorem}{Theorem}

\theoremstyle{definition}
\newtheorem{definition}{Definition}

\theoremstyle{Solution}
\newtheorem{solution}{Solution}

\begin{document}

\maketitle
\tableofcontents

\section{Vocabulary}

\begin{definition}[Theorem]
A \textit{theorem} is a statement that is true and has been proved to be true. 
\end{definition}

\begin{definition}[Proposition]
A statement that is true but not as significant as a theorem is sometimes called a \textit{proposition}. 
\end{definition}

\begin{definition}[Lemma]
A \textit{lemma} is a theorem whose main purpose is to help prove another theorem. 
\end{definition}

\begin{definition}[Corollary]
A \textit{corollary} is a result that is an immediate consequence of a theorem or proposition.
\end{definition}

\section{Basic Calculus and Analysis}

\begin{definition}[Classes of differentiability]
Let $U \subseteq \textbf{R}^n$ be open and let $f:U \rightarrow \textbf{R}^m$. A function $f$ is called continuous if

\begin{equation}
	\lim_{\textbf{y} \rightarrow \textbf{x}} f(\textbf{y}) = f(\textbf{x})
\end{equation}

for every $\textbf{x} \in U$. Such a function is said to be of class $C^0$. A function $f$ is called differentiable if its derivative exists at every point of $U$. For an integer $k \geq 1$, the function $f$ is said to be of class $C^k$ if all partial derivatives of order at most $k$ exist and are continuous on $U$. In particular, $C^1$ describes continuous first derivatives; $C^2$ describes continuous first and second derivatives, and so on. A function belonging to $C^\infty$ has continuous derivatives of every order and is called smooth.
\end{definition}

\begin{definition}[Left-hand limit]
Let $f:D_f \rightarrow \mathbb{R}$, $x \mapsto y=f(x)$, be a function and let $x_0 \in \mathbb{R}$. We say that $f(x)$ converges to $L \in \mathbb{R}$ as $x$ approaches $x_0$ from the left if, for every sequence $\{x_n\}_{n \in \mathbb{N}}$ with $x_n \in D_f$, $x_n < x_0$, and

\begin{equation}
	\lim_{n \rightarrow \infty} x_n = x_0,
\end{equation}

the sequence $\{y_n\}_{n \in \mathbf{N}}$, defined by $y_n=f(x_n)$, converges to $L$. This is written as

\begin{equation}
	\lim_{x \nearrow x_0} f(x)=L
	\qquad
	\left(\text{or } \lim_{x \rightarrow x_0^-}f(x)=L\right).
\end{equation}

The number $L$ is called the left-hand limit of $f$ at $x_0$.
\end{definition}

\begin{definition}[Right-hand limit]
Let $f:D_f \rightarrow \mathbb{R}$, $x \mapsto y=f(x)$, be a function and let $x_0 \in \mathbb{R}$. We say that $f(x)$ converges to $R \in \mathbb{R}$ as $x$ approaches $x_0$ from the right if, for every sequence $\{x_n\}_{n \in \mathbb{N}}$ with $x_n \in D_f$, $x_n > x_0$, and

\begin{equation}
	\lim_{n \rightarrow \infty} x_n = x_0,
\end{equation}

the sequence $\{y_n\}_{n \in \mathbf{N}}$, defined by $y_n=f(x_n)$, converges to $R$. This is written as

\begin{equation}
	\lim_{x \searrow x_0} f(x)=R
	\qquad
	\left(\text{or } \lim_{x \rightarrow x_0^+}f(x)=R\right).
\end{equation}

The number $R$ is called the right-hand limit of $f$ at $x_0$.
\end{definition}

\begin{definition}[Limit]
Let $f:D_f \rightarrow \mathbb{R}$, $x \mapsto y=f(x)$, be a function and let $x_0 \in \mathbb{R}$. We say that $f(x)$ converges to $y_0 \in \mathbb{R}$ as $x$ approaches $x_0$ if $f(x)$ converges to $y_0$ both from the left and from the right. In this case, one writes

\begin{equation}
	\lim_{x \rightarrow x_0}f(x)=y_0
	=
	\lim_{x \searrow x_0}f(x)
	=
	\lim_{x \nearrow x_0}f(x).
\end{equation}

The number $y_0$ is called the limit of $f$ as $x$ approaches $x_0$.
\end{definition}

\begin{definition}[Differentiable function and its derivative]
Let $f:D_f \rightarrow \mathbb{R}$, $x \mapsto y=f(x)$, be a function, let $I \subseteq D_f$ be an open interval, and let $x_0 \in I$. The function $f$ is called differentiable at $x_0$ if the limit

\begin{equation}
	\lim_{\Delta x \rightarrow 0}
	\frac{f(x_0+\Delta x)-f(x_0)}{\Delta x}
\end{equation}

exists. In this case, the limit is denoted by

\begin{equation}
	\frac{d}{dx}f(x_0)
\end{equation}

or

\begin{equation}
	f'(x_0),
\end{equation}

and is called the \textit{difference quotient} or the \textit{derivative} of the function $f$ at $x_0$.
\end{definition}

\section{Unconstrained Optimization}

\begin{definition}[Interior Point] $\mathbf{x}^*$ is an interior point of the domain of $F$ if there exists $r>0$ such that the whole ball $B_r(\mathbf{x}^*)$ is contained in the domain of $F$.
\end{definition}

\begin{definition}[Global and local maxima and minima] Let $F: U \mathbf{ \rightarrow R^1}$ be a real-valued function of $n$ variables, whose domain $U$ is a subset of $\mathbb{R}^n$.
\begin{enumerate}
	\item A point $\textbf{x}^* \in U$ is a $\textbf{max}$ of $F$ on $U$ if $F(\textbf{x}^*) \geq F(\textbf{x})$ for all $\textbf{x} \in U$.
	\item $\textbf{x}^* \in U$ is a $\textbf{strict max}$ if $\textbf{x}^*$ is a max and $F(\textbf{x}^*) \textgreater F(\textbf{x})$ for all $\textbf{x} \neq \textbf{x}^*$ in $U$.
	\item $\textbf{x}^* \in U$ is a $\textbf{local (or relative) max}$ of $F$ if there is a ball $B_r(\textbf{x}^*)$ about $\textbf{x}^*$ such that $F(\textbf{x}^*) \geq F(\textbf{x})$ for all $\textbf{x} \in B_r(\textbf{x}^*) \cap U$.
	\item $\textbf{x}^* \in U$ is a $\textbf{strict local max}$ of $F$ if there is a ball $B_r(\textbf{x}^*)$ about $\textbf{x}^*$ such that $F(\textbf{x}^*) \textgreater F(\textbf{x})$ for all $\textbf{x} \neq \textbf{x}^* \in B_r(\textbf{x}^*) \cap U$.
\end{enumerate}

If we want to emphasize that a point $\textbf{x}^*$ is a max of $F$ on the whole domain of $U$, we call $\textbf{x}^*$ a $\textbf{global max}$ or $\textbf{absolute max}$ of $F$ on $U$.
\end{definition}

\begin{theorem}
Let $F: U \mathbf{ \rightarrow R^1}$ be a $C^1$ function defined on a subset $U$ of $\mathbb{R}^n$. If $\mathbf{x}^*$ is a local max or min of $F$ in $U$ and if $\mathbf{x}^*$ is an interior point of $U$, then
 \end{theorem}

\begin{align}
\frac{\partial F}{\partial x_i}(\mathbf{x}^*)=0 \quad \text{for} \quad i=1,...,n.
\end{align}

\begin{definition}[Critical Point] The $n$-vector $\mathbf{x}^*$ is a $\textbf{critical point}$ of a function $F(x_1,...,x_n)$ if it satisfies (1).
\end{definition}

\begin{theorem} 
Let $F: U \mathbf{ \rightarrow R^1}$ be a $C^2$ function whose domain is a convex open subset $U$ of $\mathbb{R}^n$.
\begin{enumerate}
	\item The following three conditions are equivalent:
	\begin{enumerate}
		\item $F$ is a concave function of $U$; and
		\item $F(\textbf{y})-F(\textbf{x}) \leq DF(\textbf{x})(\textbf{y}-\textbf{x})$ for all $\textbf{x},  \textbf{y} \in U$: and
		\item $D^2F(\textbf{x})$ is negative semidefinite for all $\textbf{x} \in U$
	\end{enumerate}
	\item The following three conditions are equivalent:
	\begin{enumerate}
		\item $F$ is a convex function on $U$; and 
		\item $F(\textbf{y})-F(\textbf{x}) \geq DF(\textbf{x})(\textbf{y}-\textbf{x})$ for all $\textbf{x},  \textbf{y} \in U$: and
		\item $D^2F(\textbf{x})$ is positive semidefinite for all $\textbf{x} \in U$
	\end{enumerate}
	\item If $F$ is a concave function on $U$ and $DF(\textbf{x}^*)=\mathbf{0}$ for some $\textbf{x}^* \in U$, then $\textbf{x}^*$ is a $\textit{global max}$ of $F$ on $U$.
	\item If $F$ is a convex function on $U$ and $DF(\textbf{x}^*)=\mathbf{0}$ for some $\textbf{x}^* \in U$, then $\textbf{x}^*$ is a $\textit{global min}$ of $F$ on $U$.
\end{enumerate}
\end{theorem}

\begin{definition}[Convex Set]
A set $u$ is a $\textbf{convex set}$ if whenever $\textbf{x}$ and $\textbf{y}$ are points in $U$, the line segment joining $\textbf{x}$ and $\textbf{y}$,

\begin{equation}
	\ell(\textbf{x,y}) \equiv {(t\textbf{x}+(1-t)\textbf{y}:0 \leq t \leq 1},
\end{equation}

is also in $U$. 
\end{definition}

\begin{definition}
(Definite, semidefinite): Positive definite means $x \textgreater 0$; positive semidefinite means $x \geq 0$.
\end{definition}

\section{Constrained Optimization}

\begin{theorem} Let $f$ and $h$ be $C^1$ functions of two variables. Suppose that $\textbf{x}^*=(x_1^*,x_2^*)$ is a solution of the problem

\begin{equation}
	\begin{aligned}
		\max_{x_1,x_2}f(x_1,x_2), \quad \text{subject to} \quad h(x_1,x_2)=c
	\end{aligned}
\end{equation}

Suppose further that $(x_1^*,x_2^*)$ is not a critical point of $h$. Then, there is a real number $\mu^*$ such that $(x_1^*,x_2^*,\mu^*)$ is a critical point of the Langrangian function

\begin{equation}
	\begin{aligned}
		L(x_1,x_2,\mu) \equiv f(x_1,x_2)-\mu[h(x_1,x_2)-c].
	\end{aligned}
\end{equation}

In other words, at $(x_1^*,x_2^*,\mu^*)$

\begin{equation}
	\begin{aligned}
		\frac{\partial L}{\partial x_1}=0, \quad \frac{\partial L}{\partial x_2}=0, \quad \frac{\partial L}{\partial \mu}=0.
	\end{aligned}
\end{equation}
\end{theorem}

\begin{definition}[Binding constraint]
If the maximizer $\textbf{p}$ lies on the boundary of the constraint set $h(x_1,x_2)$, we say the constraint is binding (or is active, effective, or tight) at $\textbf{p}$. 
\end{definition}

\begin{definition}[Not binding constraint] If the maximizer is in the interior of the constraint set, we say that the constraint is not binding (inactive, ineffective, loose) at $\textbf{p}$.
\end{definition}

\section{Homogenous and Homothetic Functions}

\begin{definition}[Homogenous of degree k] For any scalar $k$, a real-valued function $f(x_1,...,x_n)$ is homogenous of degree $k$ if
\begin{equation}
	 f(tx_1,...,tx_n)=t^kf(x_1,...,x_n) \quad \text{for all} \quad x_1,...,x_n \quad \text{and all} \quad t \textgreater 0.
\end{equation}
\end{definition}

\begin{theorem} 
Let $z=f(\textbf{x})$ be a $C^1$ function on an open cone in $\mathbb{R}^n$. If $f$ is homogenous of degree $k$, its first order partial derivatives are homogenous of degree $k-1$.
\end{theorem}

\begin{definition}[Cone]
A cone is a set with the property that whenever $\textbf{x}$ is in the set, every positive scalar multiple $t\textbf{x}$ of $\textbf{x}$ is in the set.
\end{definition}

\begin{definition}[Marginal Rate of Substitution]
The marginal rate of substitution (MRS) describes how many units of good $Y$ a consumer is willing to trade for one additional unit of good $X$. It is the quotient of the marginal utility functions. This means it tells us how much higher the utility of $X$ is compared to that of $Y$ 
\begin{equation}
	MRS: \quad \left.\frac{dY}{dX}\right|_{\hat X,\hat Y} = - \frac{\partial U(\hat X, \hat Y)/ \partial X}{\partial U(\hat X, \hat Y)/ \partial Y}
\end{equation}
\end{definition}

\begin{definition}[Marginal Rate of Technical Substitution]
The marginal rate of technical substitution (MRTS) is the same as MRS, but with a output-function instead of utility.
\end{definition}

\begin{definition}[Income Expansion Path]
The income expansion path is the parametrized curve $I \mapsto \textbf{x}(I)$ that indicates the bundle $(x_1,...,x_n)$ demanded for different income levels. 
\end{definition}

\begin{theorem}
 Let $U(\textbf{x})$ be a utility function on $\mathbb{R}^n_+$ that is homogenous of degree $k$. Then,
\begin{enumerate}
	\item the MRS is constant along the rays from the origin,
	\item income expansion paths are rays from the origin,
	\item the corresponding demand depends linearly on income, and
	\item the income elasticity of demand is identically 1.
\end{enumerate}
Let $q=f(\textbf{x})$ be a production function on $\mathbb{R}^n_+$ that is homogenous of degree $k$. Then, 
\begin{enumerate}
	\item the MRTS is constant along rays from the origin, and
	\item the corresponding cost function is homogenous of degree $1/k: C(q)=bq^{1/k}$.
\end{enumerate}
\end{theorem}

\begin{theorem} ($\textbf{Euler's Theorem}$) $\quad$ Let $f(\textbf{x})$ be a $C^1$ homogenous function of degree $k$ on $\mathbb{R}^n_+$. Then, for all $\textbf{x}$, 

\begin{equation}
	x_1\frac{\partial f}{\partial x_1}(\textbf{x})+x_2\frac{\partial f}{\partial x_2}(\textbf{x})+...+x_n\frac{\partial f}{\partial x_n}(\textbf{x})=kf(\textbf{x}),
\end{equation}

or, in gradient notation,

\begin{equation}
	\textbf{x} \cdot \nabla f(\textbf{x})=kf(\textbf{x}),
\end{equation}

or, similarly,

\begin{equation}
	\begin{pmatrix}
		x_1\\
		\vdots\\
		x_n
	\end{pmatrix}
	\begin{pmatrix}
		F_{x1}(\textbf{x})\\
		\vdots\\
		F_{xn}(\textbf{x})
	\end{pmatrix}
	=
	kF(x_1,\ldots,x_n).
\end{equation}
\end{theorem}

\begin{theorem}[Wicksell's Law]
If \(f\) is homogeneous of degree 1 and has diminishing marginal product in \(x_1\), then applying Euler to \(f_{x1}\) yields

\[
0=x_1f_{x1x1}+x_2f_{x2x1} 
\]

\[
\Rightarrow f_{x2x1}=-\frac{x_1f_{x1x1}}{x_2} > 0
\]

Hence the marginal product of one factor increases when the other factor increases.
\end{theorem}


\begin{theorem}[Homogenizing a function]
Let $(x_1,...,x_n)\mapsto f(x_1,...,x_n)$ be a real-valued function defined on a cone $C$ in \(\mathbb{R}^n\). Let \(k\) be an integer. Define a new function $F$ of $n+1$ variables by 

\begin{equation}
	F(x_1,...,x_n,z)=z^k\cdot f \left(\frac{x_1}{z},...,\frac{x_n}{z}\right).
\end{equation}

Then, $F$ is a homogenous function of degree $k$ on the cone $C \times \mathbb{R}_+$ in $\mathbb{R}^{n+1}$. Since $f(\textbf{x})=F(\textbf{x},1)$ for all $\textbf{x} \in C$, we can consider $f$ as the restriction of $F$ to an $n$-dimensional subset of $\mathbb{R}^{n+1}$.
\end{theorem}

\vspace{1em}

Proposition 3.1 (Radial expansions): The level sets of a homogenous (production) function are radial expansions and contractions of each other.

\vspace{1em}

\begin{definition}[Monotonously decreasing function] $f$ is monotonously falling if $f(x) \geq f(x') \forall x,x' \in X$ such that $x \leq x'$.
\end{definition}

\begin{definition}[Monotonic transformation] 
Let $I$ be an interval on the real line. Then, $g:I\rightarrow \textbf{R}$ is a $\textbf{monotonic transformation}$ of $I$ if $g$ is a strictly increasing function on $I$. Furthermore, if $g$ is a monotonic transformation and $u$ is a real-valued function of $n$ variables, then we say that

\begin{equation}
	g \circ u: \mathbf{x} \mapsto g(u(\mathbf{x}))
\end{equation}

is a $\textbf{monotonic transformation of u}$. For instance, the functions $3z+2, \quad z^2, \quad e^z, \quad \text{and} \quad \text{ln} z$ are all monotonic transformations of $\mathbb{R}_{++}$. Consequently, the utility functions $3xy+2, \quad (xy)^2, \quad e^{xy}, \quad \text{and} \quad \text{ln} x+\text{ln} y$ are monotonic transformations of the utility function $u(x,y)=xy$.
\end{definition}

\begin{definition}[Ordinal and cardinal]
A characteristic of functions is called ordinal if every monotonic transformation of a function with this characteristic still has that characteristic. Cardinal properties are not preserved by monotonic transformations.
\end{definition}

\begin{definition}[Homothetic function]
A function $v: \mathbb{R}^n_+ \rightarrow \mathbb{R}$ is called $\textbf{homothetic}$ if it is a monotone transformation of a homogenous function, that is, if there is a monotonic transformation $z \mapsto g(z)$ of $\mathbb{R}_+$ and a homogenous function $u:\mathbb{R}^n_+ \rightarrow \mathbb{R}_+$ such that $v(\textbf{x})=g(u(\textbf{x}))$ for all $\textbf{x}$ in the domain.
\end{definition}

\begin{theorem} 
Let $u: \mathbb{R}^n_+ \rightarrow \textbf{R}$ be a strictly monotonic function. Then, $u$ is homothetic if and only if for all $\textbf{x}$ and $\textbf{y}$ in $\mathbb{R}^n_+$,

\begin{equation}
	u(\textbf{x}) \geq u(\textbf{y}) \quad \Leftrightarrow \quad u(\alpha \textbf{x}) \geq u(\alpha \textbf{y}) \quad \forall \quad \alpha > 0.
\end{equation}
\end{theorem}

\section{Concave and Quasiconcave Functions}

\begin{theorem}[Convex domain]
Concave and convex functions are required to have convex domains. They contain every line segment between any two points in the set:

\begin{equation}
x_1,x_2 \in C, \quad 0 \leq \theta \leq 1 \Rightarrow \theta x_1+(1-\theta)x_2 \in C.
\end{equation}
\end{theorem}

\begin{definition}[Concave and convex $n$-dimensional functions]
A real-valued function $f$ defined on a convex subset $U$ of $\mathbb{R}^n$ is $\textbf{concave}$ if for all $\textbf{x, y}$ in $U$ and for all $t$ between $0$ and $1$,

\begin{equation}
	f(t\mathbf{x}+(1-t)\textbf{y}) \geq tf(\textbf{x})+(1-t)f(\textbf{y}).
\end{equation}

A real-valued function $g$ defined on a convex subset $U$ of $\mathbb{R}^n$ is $\textbf{convex}$ if for all $\textbf{x, y}$ in $U$ and for all $t$ between $0$ and $1$,

\begin{equation}
	g(t\mathbf{x}+(1-t)\textbf{y}) \leq tg(\textbf{x})+(1-t)g(\textbf{y}).
\end{equation}
\end{definition}

\begin{theorem} 
Let $f$ be a $C^1$ function on an interval $I$ in $\textbf{R}$. Then, $f$ is concave on $I$ if and only if

\begin{equation}
	f(y)-f(x) \leq f'(x)(y-x) \quad \text{for all} x,y \in I.
\end{equation}

The function $f$ is convex on $I$ if and only if

\begin{equation}
		f(y)-f(x) \geq f'(x)(y-x) \quad \text{for all} x,y \in I.
\end{equation}
\end{theorem}

\begin{theorem}
Let $f$ be a $C^1$ function on a convex subset $U$ of $\mathbb{R}^n$. Then, $f$ is concave on $U$ if and only if for all $\textbf{x, y}$ in $U$:

\begin{equation}
	f(\textbf{y})-f(\textbf{x}) \leq Df(\textbf{x})(\textbf{y}-\textbf{x});
\end{equation}

that is,

\begin{equation}
	f(\textbf{y})-f(\textbf{x}) \leq \frac{\partial f}{\partial x_1}(\textbf{x})(y_1-x_1) +...+\frac{\partial f}{\partial x_n}(\textbf{x})(y_n-x_n).
\end{equation}

Similarly, $f$ is convex on $U$ if and only if $f(\textbf{y})-f(\textbf{x}) \geq Df(\textbf{x})(\textbf{y}-\textbf{x})$ for all $\textbf{x, y}$ in $U$.
\end{theorem}

\begin{definition}[Quasiconcavity and Quasiconvexity]
Let \(f:\mathcal{S}\mapsto\mathbb{R}\).\par
\begin{enumerate}
	\item The function \(f\) is quasiconcave on \(\mathcal{S}\) if and only if the upper level set \(C_a^+=\{x\in\mathcal{S}:f(x)\geq a\}\) is a convex set for every \(a\in\mathcal{S}\).
	\item The function \(f\) is quasiconvex on \(\mathcal{S}\) if and only if the lower level set \(C_a^-=\{x\in\mathcal{S}:f(x)\leq a\}\) is a convex set for every \(a\in\mathcal{S}\).
\end{enumerate}
\end{definition}

\section{Difference Equations}

\begin{definition}[Functional Equation]
A \textit{functional equation} is an equation where the unknown is the function.
\end{definition}

\begin{definition}[Ordinary difference equation]
An \textit{ordinary difference equation} is a functional equation involving one or more of the differences $\Delta y, \Delta^2y$, etc. of a function of time.
\end{definition}

\begin{definition}[Linear, constant-coefficient non-homogenous difference equation]
The general $n$-th order form of such equations is
\begin{equation}
	c_ny_{t+n}+c_{n-1}y_{t+n-1}+...+c_{1}y_{t+1}+c_{0}y_{t}=g(t),
\end{equation}
where the $c$'s are given constants and $g(t)$ is a known function.
\end{definition}

\begin{definition}[Linear, constant-coefficient homogenous difference equation]
The general $n$-th order form of such equations is
\begin{equation}
	c_ny_{t+n}+c_{n-1}y_{t+n-1}+...+c_{1}y_{t+1}+c_{0}y_{t}=0.
\end{equation}
\end{definition}

\begin{theorem}[]
If $y_1(t)$ is a solution of the homogenous equation, then $Ay_1(t)$, where $A$ is an arbitrary constant, is also a solution.
\end{theorem}

\begin{theorem}[]
If $y_1(t), y_2(t)$ are two distinct (linearly independent) solutions of the homogenous equation $(n>1)$, then $A_1y_1(t)+A_2y_2(t)$ is also a solution for any two constants $A_1, A_2$.
\end{theorem}

\begin{theorem}[]
If $\overline{y}(t)$ is a particular solution of the non-homogenous equation, the general solution of this same equation is obtained adding $\overline{y}(t)$ to the general solution of the corresponding homogenous equation.
\end{theorem}

\begin{theorem}[General form of first-order non-homogenous equations]
The general form of these equations is
\begin{equation}
	c_1y_t+c_0y_{t-1}=g(t),
\end{equation}
where $c_0,c_1$ are given in constants and $g(t)$ is a known function and its general solution is $y_t=A(-b)^t+\overline{y}_t \Leftrightarrow A=\frac{y^*-\overline{y}_{t^*}}{(-b)^{t^*}}$.
\end{theorem}

\begin{theorem}[General form of first-order homogenous equations]
The general form of these equations is
\begin{equation}
	c_1y_t+c_0y_{t-1}=0,
\end{equation}
or
\begin{equation}
	y_t+by_{t-1}=0,
\end{equation}
where $b\equiv c_0/c_1$ and the solution is $y_t=A(-b)^t \Leftrightarrow A=y^*/(-b)^{t*}.$
\end{theorem}

\begin{theorem}
(Monotony and convergence characteristics of the general solution of first-order difference equations with constant coefficients):
There are only four possible outcomes: Monotonous and convergent $A \in (0,1)$, oscillating and convergent $A \in (-1,0)$, monotonous and divergent $A \in (1,\infty)$, and oscillating and divergent $A \in (-\infty,-1]$.
\end{theorem}

\begin{theorem}[First-order systems]
The following is a first order non-homogenous system

\begin{equation}
\begin{aligned}
	y_{t+1}&=a_{11}y_{t}+a_{12}z_{t}+g_{1}(t), \\
	z_{t+1}&=a_{21}y_{t}+a_{22}z_{t}+g_{2}(t),
\end{aligned}
\end{equation}

where the coefficients $a_{ij}$ are given constants and $g_1(t),g_2(t)$ are known functions. The corresponding homogenous system is

\begin{equation}
\begin{aligned}
	y_{t+1}&=a_{11}y_{t}+a_{12}z_{t}, \\
	z_{t+1}&=a_{21}y_{t}+a_{22}z_{t}.
\end{aligned}
\end{equation}

We solve (37) by solving the second equation for $z_t$, creating a duplicate where time is shifted by one period, taking the first equation in (37) and shifting it too and then substituting the $z$-terms with the ones we just derived. This eventually gives us the second-order difference equation

\begin{equation}
	y_{t}-(a_{11}+a_{22})y_{t-1}-(a_{12}a_{21}-a_{11}a_{22})y_{t-2}=0
\end{equation}
with only one unknown $y_t$.
\end{theorem}

\begin{solution}[General solution of the homogenous system: Direct method]
We solve the system from equation (36) directly by taking it's homogenous counterpart and then substituting $y_t=\alpha_1\lambda^t, z_t=\alpha_2\lambda^t$. This gives us
\begin{equation}
	\begin{aligned}
		\alpha_1\lambda^{t+1}&=a_{11}\alpha_1\lambda^t+a_{12}\alpha_2\lambda^t, \\
		\alpha_2\lambda^{t+1}&=a_{21}\alpha_1\lambda^t+a_{22}\alpha_2\lambda^t.
	\end{aligned}
\end{equation}
This is equivalent to
\begin{equation}
	\begin{aligned}
		\lambda^t[\alpha_1\lambda-a_{11}\alpha_1-a_{12}\alpha_2]&=0, \\
		\lambda^t[\alpha_2\lambda-a_{21}\alpha_1-a_{22}\alpha_2]&=0,
	\end{aligned}
\end{equation}
leading us to
\begin{equation}
	\begin{aligned}
		\alpha_1(\lambda-a_{11})-a_{12}\alpha_2&=0, \\
		\alpha_2(\lambda-a_{22})-a_{21}\alpha_1&=0.
	\end{aligned}
\end{equation}
We must invert the sign and take the determinant of the matrix $A$ which represents that consequential system. Remember that the variables $\alpha_1, \alpha_2$ must be dropped out of the equation.
\begin{equation}
	\begin{aligned}
		det(A)&=\left|\begin{matrix}
			(a_{11}-\lambda) & a_{12} \\ 
			a_{21} & (a_{22}-\lambda) 
		\end{matrix} \right| \\
		&\Leftrightarrow (a_{11}-\lambda)(a_{22}-\lambda)-a_{21}a_{12}=0 \\
		&\Leftrightarrow a_{11}a_{22}-a_{11}\lambda-\alpha_{22}\lambda+\lambda^2-a_{21}a_{12}=0 \\
		&\Leftrightarrow \lambda^2-(a_{11}+a_{22})\lambda+(a_{11}a_{22}-a_{21}a_{12})=0.
	\end{aligned}
\end{equation}
The determinantal equation in the first equality and its expanded form are called the characteristic equation of the system of difference equations. From Theorem 18 we know that the constants $a_{ij}$ are known, such that the last equation will give us two solutions $\lambda_1,\lambda_2$.
\end{solution}

\section{Differential Equations}

\begin{theorem}[General form of first-order non-homogenous equations]
The general form of these equations is
\begin{equation}
	\begin{aligned}
		\dot{x}(t)=a_{11}x(t)+a_{12}y(t)+g(t), \\
		\dot{y}(t)=a_{21}x(t)+a_{22}y(t)+g(t),
	\end{aligned}
\end{equation}
where $a_{ij}$ are given constants and $g(t)$ is a known function.
\end{theorem}

\begin{solution}[General solution of the homogenous system: Direct method]
We solve the system from equation (43) directly by taking it's homogenous counterpart and setting $x(t)=\alpha_1e^{\lambda t}, y(t)=\alpha_2e^{\lambda t}$. This gives us
\begin{equation}
	\begin{aligned}
		\alpha_1\lambda e^{\lambda t}&=a_{11}\alpha_1e^{\lambda t}+a_{12}\alpha_2e^{\lambda t}, \\
		\alpha_2\lambda e^{\lambda t}&=a_{21}\alpha_1e^{\lambda t}+a_{22}\alpha_2e^{\lambda t}.
	\end{aligned}
\end{equation}

This is equivalent to

\begin{equation}
	\begin{aligned}
		e^{\lambda t}[(a_{11}-\lambda)\alpha_1+a_{12}\alpha_2]&=0, \\
		e^{\lambda t}[a_{21}\alpha_1+(a_{22}-\lambda)\alpha_2]&=0.
	\end{aligned}
\end{equation}
leading us to
\begin{equation}
	\begin{aligned}
		(a_{11}-\lambda)\alpha_1+a_{12}\alpha_2&=0, \\
		a_{21}\alpha_1+(a_{22}-\lambda)\alpha_2&=0.
	\end{aligned}
\end{equation}
We must take the determinant of the matrix $A$ which represents that consequential system. Remember that the variables $\alpha_1, \alpha_2$ must be dropped out of the equation.
\begin{equation}
	\begin{aligned}
		det(A)&=\left|\begin{matrix}
			(a_{11}-\lambda) & a_{12} \\ 
			a_{21} & (a_{22}-\lambda) 
		\end{matrix} \right| \\
		&\Leftrightarrow (a_{11}-\lambda)(a_{22}-\lambda)-a_{21}a_{12}=0 \\
		&\Leftrightarrow a_{11}a_{22}-a_{11}\lambda-\alpha_{22}\lambda+\lambda^2-a_{21}a_{12}=0 \\
		&\Leftrightarrow \lambda^2-(a_{11}+a_{22})\lambda+(a_{11}a_{22}-a_{21}a_{12})=0.
	\end{aligned}
\end{equation}
The determinantal equation in the first equality and its expanded form are called the characteristic equation of the system of difference equations. From Theorem 18 we know that the constants $a_{ij}$ are known, such that the last equation will give us two solutions $\lambda_1,\lambda_2$.
\end{solution}

\section{Optimal Control}


\end{document}